\section{Every even supersingular thin system}\label{sec:ss-all-even}

We now give a single section for every even first degree. Its Cartier kernel is a fixed rank-three bundle, and two residue functionals show that the image is already stable. This removes any divisibility-by-three condition from the supersingular thin theorem.

\begin{proposition}[An explicit section in every even degree]\label{prop:ss-all-even}
On the model $E/\FF_2$ given by $y^2+y=x^3$, write $z,w$ for the coordinates on a second copy. Let $m=2t\ge6$, and put
\[
 f=\begin{cases}
 x^{t/2},&t\text{ even},\\
 x^{(t-3)/2}y,&t\text{ odd},
 \end{cases}
 \qquad
 a=x+f^2,\qquad c=x^{t-2}y.
\]
The section
\[
                         s=a+z+cw\in H^0(E\times E,N(m,3))
\]
has first and stable Cartier rank $3m-2$. After extension to any algebraically closed field of characteristic two, a general smooth member of $|N(m,3)|$ has genus $3m+1$ and full and relative $2$-rank $3m-2$.
\end{proposition}
\begin{proof}
We carry out the bundle and residue calculations over $\FF_2$; all the displayed functions are defined over this field. Write $\omega=dx$ and $C(h)=C_E(h\omega)/\omega$. Since $x$ is a $2$-basis,
\[
 C(u^2+xv^2)=v,\qquad
 \partial(u^2+xv^2)=v^2.
\]
In the second factor,
\[
 C'(1)=C'(z^2)=C'(w^2)=0,\quad
 C'(z)=1,\quad C'(w)=z,\quad C'(zw)=w.
\]
These identities follow by substituting $w^2+w=z^3$ into the same coefficient rule. Hence, identifying a section of $N(m,3)$ with a triple in $H^0(E,\OO_E(mO))^3$, multiplication followed by Cartier is
\begin{equation}\label{eq:ss-even-cartier}
 T_s(t_0,t_1,t_2)=
 \bigl(C(t_0+at_1),\ C(ct_0+at_2),\ C(ct_1+t_2)\bigr).
\end{equation}

Introduce
\[
 (g,h)=
 \begin{cases}
 (x^{t/2},\,x^{(t-2)/2}y),&t\text{ even},\\
 (x^{(t-3)/2}y,\,x^{(t+1)/2}),&t\text{ odd},
 \end{cases}
 \qquad H=fg+h.
\]
The elliptic equation gives
\begin{equation}\label{eq:ss-even-identities}
 c=h^2+xg^2,\qquad \partial a=1,\qquad
 \partial c=g^2,\qquad ag^2+c=H^2.
\end{equation}
The functions $f,g,h$ have respective pole orders $t,t,t+1$ at $O$, with leading coefficient one for the parameter $\tau=x/y$. Thus $a,c,H$ have respective pole orders $m,m-1,m$, and $s$ is a section of the asserted bundle.

\paragraph{The exact kernel bundle.}
Put $L_1=\OO_E(mO)$, and let $F:E\to E$ be Frobenius on the $\FF_2$ model. Formula~\eqref{eq:ss-even-cartier} defines an $\OO_E$-linear map
\[
                   T:(F_*L_1)^3\longrightarrow L_1^3.
\]
The source module structure is through squares. We claim that
\begin{equation}\label{eq:ss-even-kernel}
 0\longrightarrow F_*\OO_E\oplus\OO_E
 \xrightarrow{\ i\ }(F_*L_1)^3
 \xrightarrow{\ T\ }L_1^3\longrightarrow0,
 \qquad i(q,r)=qs+r^2\partial s
\end{equation}
is exact, where $\partial s=(1,0,g^2)$.

Over the function field, put $q=t_1$ for a vector in the kernel and subtract $qs$. The first and third equations in~\eqref{eq:ss-even-cartier} leave a vector $(u,0,v)$ with $C(u)=C(v)=0$, so $u=r^2$ and $v=r_2^2$. The middle equation is $rg+r_2=0$; hence the vector is exactly $qs+r^2\partial s$. Conversely, substitution proves that every such vector is in the kernel.

Away from $O$, regularity of the vector gives regularity of $q=t_1$ and then of $r^2=t_0+aq$, hence of $r$. At $O$, the coefficient matrix for $q,r^2$ in the first and third components, after multiplying them by $\tau^m$, is
\[
 \begin{pmatrix}
 \tau^ma&\tau^m\\
 \tau^mc&\tau^mg^2
 \end{pmatrix}.
\]
Its determinant is $\tau^{2m}H^2$, a unit. Thus membership in $L_1^3$ at $O$ is equivalent to regularity of $q,r$. This proves the kernel assertion as a statement of bundles.

The quotient by that kernel has rank three and degree $3m$: finite Frobenius pushforward preserves Euler characteristic, and on an elliptic curve Euler characteristic equals degree. The target $L_1^3$ also has degree $3m$. The induced generically injective map between these bundles consequently has torsion cokernel of length zero. This proves the surjectivity in~\eqref{eq:ss-even-kernel}.
Since $H^1(E,L_1)=0$, global sections give
\begin{equation}\label{eq:ss-even-global-kernel}
 \ker T_s=\langle s,\partial s\rangle,\qquad
 \operatorname{coker}T_s
       \simeq H^1(E,F_*\OO_E\oplus\OO_E).
\end{equation}
Both spaces have dimension two, so $\rk T_s=3m-2$.

\paragraph{The first residue functional.}
For an output triple $u=(u_0,u_1,u_2)$, set
\[
 \ell(u)=\operatorname{Res}_O
                \frac{gu_0+u_2}{H}\,\omega.
\]
It kills the image of $T_s$. Indeed, by~\eqref{eq:ss-even-identities},
\[
 \frac{g}{H}T_0(t)+\frac1H T_2(t)
        =C(t_1)+C\!\left(\frac{g^2t_0+t_2}{H^2}\right).
\]
The first differential has no possible pole away from $O$, so its residue there is zero by the global residue theorem. The argument of the second Cartier operator is regular at $O$, because its numerator has pole order at most $2m$ and its denominator has pole order $2m$. Cartier preserves regular differentials.

Put
\[
                       R=\frac{c+xg+fh}{H}.
\]
The terms $c/H$ and $xg/H$ are regular at $O$. The term $fh/H$ has a simple pole there with coefficient one. Since $\omega=(1+O(\tau))\,d\tau$, this gives
\begin{equation}\label{eq:ss-even-residue}
                   \operatorname{Res}_O(R\omega)=1.
\end{equation}
The identity $(ag+c)/H=f+R$ and the global residue theorem for $f\omega$ imply $\ell(s)=1$. On the other hand, $(g+g^2)/H$ is regular at $O$, so $\ell(\partial s)=0$.

\paragraph{The second residue functional.}
Project the connecting map from~\eqref{eq:ss-even-global-kernel} onto the second factor $H^1(E,\OO_E)$ and pair with $\omega$ by residue duality. Denote the resulting cokernel functional by $\eta$. The residue of a principal part times $\omega$ is well-defined: regular terms and rational functions with their only pole at $O$ have zero residue. It is nonzero on $\tau^{-1}$, and hence identifies the one-dimensional group $H^1(E,\OO_E)$ with $\FF_2$.

We calculate $\eta(\partial s)$ by local preimages. On $E\setminus\{O\}$ take
\[
                  t_U=(x,0,h^2+ag^4).
\]
Formula~\eqref{eq:ss-even-cartier} gives $T_s(t_U)=\partial s$: use $C(x)=C(a)=1$, $C(h^2)=0$, and
\[
 C(cx)=h,\qquad
 C(ah^2+a^2g^4)=h.
\]
Near $O$ put
\[
 q=\frac{ag^4+c}{H^2},\qquad r=H+R,\qquad t_O=(0,q,0).
\]
The pole order of $q$ is at most $m$. After excluding the other zeros of $H$, the vector $t_O$ is therefore a local section of $L_1^3$. Direct substitution in the definitions gives
\begin{equation}\label{eq:ss-even-square}
                              r^2=x+aq.
\end{equation}
To verify this without an implicit cancellation, multiply by $H^2$. The two sides become equal because
\[
 H^4+c^2+x^2g^2+f^2h^2
   =xH^2+a^2g^4+ac
   =f^4g^4+x^2g^4+x^2g^2+f^2h^2.
\]
Furthermore,
\[
 qc+r^2g^2=q(c+ag^2)+xg^2
           =ag^4+c+xg^2=ag^4+h^2.
\]
Consequently $t_U+t_O=qs+r^2\partial s$, so $t_O$ is also a preimage of $\partial s$ and the connecting cocycle is represented by $(q,r)$. Its second residue is
\[
 \eta(\partial s)=\operatorname{Res}_O(r\omega)
 =\operatorname{Res}_O(H\omega)+\operatorname{Res}_O(R\omega)=1.
\]
Here $H$ has no pole away from $O$, and~\eqref{eq:ss-even-residue} supplies the second term.

The two cokernel functionals on the kernel basis $s,\partial s$ therefore have matrix
\[
                         \begin{pmatrix}1&0\\ *&1\end{pmatrix}.
\]
It is invertible. Kernel and image intersect trivially and have complementary dimensions, so the image is stable and the first rank already equals stable rank.

Finally, all coefficients of the section and its Cartier matrix are in $\FF_2$. After extension to $k$, absolute Cartier acts by that same matrix followed by inverse Frobenius on coefficients. The stable rank is therefore unchanged; no untwisted matrix power has replaced an arbitrary semilinear iterate. The resulting nonempty iterated-rank open meets the smooth open of the very ample complete system. Lemmas~\ref{lem:cohom} and~\ref{lem:upper} give connectedness, genus, and the matching bound $3m-2$. This proves the proposition.
\end{proof}
